\documentclass[11pt]{article}

\usepackage[T1]{fontenc}
\usepackage[utf8]{inputenc}
\usepackage{lmodern}
\usepackage[french]{babel}
\usepackage[a4paper,margin=2.6cm]{geometry}
\usepackage{amsmath}
\usepackage{amssymb}
\usepackage{booktabs}
\usepackage{array}
\usepackage{microtype}
\usepackage{listings}
\usepackage{xcolor}
\usepackage[htt]{hyphenat}
\usepackage[colorlinks=true,linkcolor=black,citecolor=black,urlcolor=blue!50!black]{hyperref}
\usepackage{xurl}

\definecolor{codegray}{gray}{0.45}
\definecolor{codebg}{gray}{0.97}
\lstset{
  basicstyle=\ttfamily\small,
  backgroundcolor=\color{codebg},
  commentstyle=\color{codegray},
  keywordstyle=\bfseries,
  breaklines=true,
  columns=fullflexible,
  keepspaces=true,
  frame=single,
  framerule=0.2pt,
  xleftmargin=1em,
  xrightmargin=1em,
}

\title{Contrôle d'admission stigmergique et de quorum~: guide du praticien}
\author{\texttt{moderation-pipeline-prototype}\thanks{Complément à
    l'article \emph{Une porte d'admission à décroissance pondérée, à
    escompte d'indépendance des sources et à seuil adaptatif de charge,
    pour un étage d'argumentation NP-complet} (même dossier), qui porte
    le modèle formel, les deux propositions et leurs preuves, le
    corollaire sur la tension d'étalonnage, ainsi que la recherche
    d'antériorités. Ce guide porte la matière appliquée~: parcours
    commenté, procédure de réglage, dépannage, adaptation.}}
\date{Août 2026}

\begin{document}
\maketitle

\section{Quand ce guide s'applique}
\label{sec:when}

Les trois conditions suivantes, réunies~:

\begin{enumerate}
\item Votre système possède une voie rapide et peu coûteuse, ainsi qu'un
  recours coûteux, dans le pire cas exponentiel (ou autrement onéreux),
  vers lequel basculent les dossiers dès que la voie rapide ne tranche
  pas.
\item L'admission au recours coûteux est aujourd'hui subordonnée à peu de
  chose~: dès qu'un dossier reçoit un premier signal de confirmation
  plausible, il passe.
\item Un adversaire peut fabriquer ce premier signal de confirmation à
  bas coût, en volume, en une seule salve, sans avoir à soutenir quoi que
  ce soit ni à coordonner des identités dans la durée.
\end{enumerate}

Le déploiement de référence est un pipeline de modération de contenu à
quatre étages, dont l'étage~2b bascule vers l'énumération, NP-complète,
des extensions préférées d'un cadre d'argumentation abstraite (Dung,
1995) lorsque l'extension fondée est vide et que le délai d'escalade
d'un dossier a expiré. Rien, ci-dessous, n'est propre aux cadres
d'argumentation~: la porte traite l'étage coûteux comme un réservoir de
ressources opaque, muni d'un signal de libération, et s'insère à toute
frontière d'admission correspondant aux trois conditions ci-dessus (files
d'attente qui se rabattent sur un solveur exhaustif, escalades de
modération vers une révision humaine, accès limité en débit à un appel de
modèle coûteux).

\medskip
\noindent
\textbf{Ce que l'on obtient.} Une salve de taille quelconque, en un seul
pas de temps, issue d'un nombre quelconque d'identités distinctes ou
partagées, est bloquée sans condition. Un attaquant patient, soutenu et
diversifié en identité n'est pas bloqué, mais son admission est retardée
d'un facteur d'environ cinq sous le réglage de référence --- la
section~\ref{sec:tuning} montre que ce facteur est un cadran, non une
constante, et ce qu'il en coûte de le tourner.

\section{La porte en trois règles}
\label{sec:rules}

\begin{enumerate}
\item \textbf{Décroître à chaque pas, avant toute autre chose.} Chaque
  dossier porte un score qui diminue d'une fraction fixe $\rho$ à chaque
  pas de temps discret, qu'il ait reçu ou non un renforcement à ce pas.
  C'est la règle ordinaire d'évaporation des phéromones
  ($\tau \leftarrow (1-\rho)\tau$) tirée de l'optimisation par colonies
  de fourmis, appliquée ici à un score de pression d'admission plutôt
  qu'à la force d'une piste. La décroissance s'exécute \emph{en premier},
  de sorte qu'un dossier qui devient silencieux commence à perdre du
  terrain immédiatement, sans période de grâce.
\item \textbf{Plafonner la contribution du pas, et escompter les
  répétitions d'une même source.} À l'intérieur d'un pas, dédupliquer
  d'abord les sources qui renforcent (une source qui soumet cinq
  éléments en un pas constitue un seul événement de renforcement, non
  cinq), puis sommer~: 1,0 par source vue pour la première fois au cours
  des $w$ derniers pas, une petite fraction $\delta$ par source revue à
  l'intérieur de cette fenêtre. Plafonner la somme à un montant fixe
  $\gamma$ \emph{après} l'étape de déduplication et d'escompte, jamais
  avant. Inverser cet ordre est l'erreur d'implémentation la plus
  lourde de conséquences de toute cette porte --- la
  section~\ref{sec:pieges} la traite en premier, précisément parce
  qu'elle s'est produite ici.
\item \textbf{Relever la barre quand l'étage coûteux est occupé.} Le
  seuil d'admission n'est pas une constante~; c'est
  $\theta_0(1 + \kappa L/C)$, où $L$ est le nombre de dossiers occupant
  l'étage coûteux, à l'échelle du système, à l'instant présent --- non un
  compte cumulatif de tout ce qui a jamais été admis. Se tromper sur
  cette distinction est le second bogue de la section~\ref{sec:pieges}.
\end{enumerate}

\section{Parcours commenté}
\label{sec:parcours}

L'implémentation de référence est \texttt{src/admission\_control.py},
classe \texttt{StigmergicQuorumAdmissionControl}. Ce qui suit reprend ce
fichier dans l'ordre de lecture, en déplaçant le raisonnement des
commentaires de code vers la prose.

\subsection{Les paramètres réglables et leurs valeurs par défaut}

\begin{lstlisting}
DEFAULT_DECAY_RATE = 0.10           # rho
DEFAULT_TICK_CAP = 1.0              # gamma
DEFAULT_INDEPENDENCE_WINDOW = 3     # w, en pas
DEFAULT_REPEAT_DISCOUNT = 0.05      # delta
DEFAULT_BASE_THRESHOLD = 4.0        # theta_0
DEFAULT_LOAD_SENSITIVITY = 0.5      # kappa
DEFAULT_CAPACITY = 5                # C
\end{lstlisting}

Ces sept nombres ne sont pas des cadrans indépendants à régler un par un
au jugé. Le corollaire~1 de l'article prouve que le rapport
$x = \theta_0\rho/\gamma$, évalué à charge nulle, gouverne \emph{à la
fois} la facilité d'admission d'un litige authentique et modérément
corroboré \emph{et} le nombre de pas soutenus dont un attaquant
pleinement saturé a besoin. Il n'existe aucune façon de déplacer $x$ qui
avantage un côté sans désavantager l'autre. Les valeurs de référence
ci-dessus ont été choisies en mesurant directement ce compromis (non en
les réglant à rebours d'un résultat de test souhaité)~: elles exigent
d'un attaquant environ cinq pas d'effort à saturation complète, tout en
gardant la barre franchissable pour un litige recevant un renforcement
indépendant à un taux de l'ordre de 40 à 50\,\% des pas en moyenne. Si le
taux de renforcement des litiges authentiques diffère dans votre
déploiement, redérivez $x$ pour votre propre population avant de toucher
à une seule constante (section~\ref{sec:tuning}).

\subsection{L'état par dossier et la boucle de pas}

\begin{lstlisting}
@dataclass
class CaseAdmissionState:
    score: float = 0.0
    last_reinforced_tick: dict[str, int] = field(default_factory=dict)
    admitted: bool = False
    admitted_tick: int | None = None
\end{lstlisting}

\texttt{last\_reinforced\_tick} est indexé par identifiant de cohorte, non
par événement~: c'est toute la mémoire que le mécanisme garde de qui a
récemment renforcé ce dossier, et c'est ce qui rend possible l'escompte
d'indépendance sans conserver l'historique complet des événements.

\begin{lstlisting}
def tick(self, t: int, reinforcements: dict[str, list[str]]) -> set[str]:
    newly_admitted: set[str] = set()
    touched = set(self._cases) | set(reinforcements)
    for case_id in touched:
        state = self._cases.setdefault(case_id, CaseAdmissionState())
        if state.admitted:
            continue
        # (1) decroissance d'abord
        state.score *= (1 - self.decay_rate)
        # (2) accumuler la contribution du pas, dedupliquee,
        #     ponderee, plafonnee
        contrib = 0.0
        for cohort in set(reinforcements.get(case_id, [])):
            last = state.last_reinforced_tick.get(cohort)
            fresh = last is None or (t - last) >= self.independence_window
            contrib += 1.0 if fresh else self.repeat_discount
            state.last_reinforced_tick[cohort] = t
        contrib = min(contrib, self.tick_cap)
        state.score += contrib
        # (3) admettre seulement si le seuil adaptatif est franchi
        if state.score >= self.current_threshold():
            state.admitted = True
            state.admitted_tick = t
            self._pending_admitted.add(case_id)
            newly_admitted.add(case_id)
    return newly_admitted
\end{lstlisting}

Notez \texttt{set(reinforcements.get(case\_id, []))} dans la boucle sur
les cohortes --- cet appel à \texttt{set()} est l'étape de déduplication,
et il s'exécute avant la logique d'escompte, non après. C'est le correctif
du premier bogue de la section~\ref{sec:pieges}~; si vous adaptez cette
porte à un autre langage ou réécrivez la boucle, placez la déduplication
là, et nulle part plus loin.

\begin{lstlisting}
def current_threshold(self) -> float:
    load = len(self._pending_admitted)
    return self.base_threshold * (1 + self.load_sensitivity * load / self.capacity)

def release(self, case_id: str) -> None:
    self._pending_admitted.discard(case_id)
\end{lstlisting}

\texttt{release()} n'est pas une plomberie facultative que l'on peut
ignorer lors d'une première intégration. Le second bogue de la
section~\ref{sec:pieges} montre ce qui arrive si on l'omet.

\subsection{L'auxiliaire en forme close}

\begin{lstlisting}
def min_ticks_to_admission(threshold, tick_cap, decay_rate) -> float:
    ceiling_fraction = threshold * decay_rate / tick_cap
    if ceiling_fraction >= 1.0:
        return math.inf
    return math.log(1 - ceiling_fraction) / math.log(1 - decay_rate)
\end{lstlisting}

Cette fonction n'est pas un utilitaire annexe~; c'est la proposition~2 de
l'article compagnon rendue exécutable, et elle sert deux fois~: une fois
pour \emph{concevoir} \texttt{base\_threshold} en fonction d'un coût
attaquant visé, et une fois à l'intérieur de la suite de tests
adversariaux, pour vérifier que la porte simulée admet effectivement un
attaquant saturé au pas prédit par la forme close. Si vous modifiez
\texttt{decay\_rate}, \texttt{tick\_cap} ou \texttt{base\_threshold},
relancez cette fonction avant de relancer la simulation --- si les deux
divergent de plus d'un pas, quelque chose dans l'implémentation s'est
éloigné du modèle qu'elle est censée réaliser.

\section{Procédure de réglage}
\label{sec:tuning}

Ne réglez pas les sept constantes indépendamment. Suivez cet ordre~:

\begin{enumerate}
\item \textbf{Fixez $\rho$ à partir de la granularité de vos pas.} $\rho$
  fixe la demi-vie de la pression non renforcée, $\ln 2/\rho$ pas.
  Choisissez d'abord une durée de pas correspondant à la fréquence à
  laquelle votre système peut plausiblement observer des événements de
  renforcement (des minutes, non des millisecondes, pour un pipeline de
  modération humaine), puis choisissez $\rho$ pour que la demi-vie
  corresponde à la durée pendant laquelle un litige réellement urgent
  devrait rester « chaud » entre deux renforcements indépendants.
\item \textbf{Fixez $\gamma$ et $\delta$ à partir de votre plafond de
  diversité d'identité.} $\gamma=1$ (une seule source fraîche par pas
  sature entièrement le plafond) est la valeur par défaut prudente.
  $\delta$ doit être assez petit pour que le taux de répétition attendu
  d'une même cohorte ne puisse atteindre le seuil même sur un temps
  illimité --- vérifiez-le avec la proposition~1~: le plafond vaut
  $\delta\gamma/\rho$, et il doit rester sous $\theta_0$ avec une marge
  qui vous convient, non tout juste en dessous.
\item \textbf{Fixez $\theta_0$ à partir de la tension d'étalonnage,
  explicitement.} Calculez $x = \theta_0\rho/\gamma$ pour plusieurs
  valeurs candidates de $\theta_0$ et, pour chacune, calculez à la fois
  (a) la probabilité qu'un litige authentique et modérément corroboré de
  votre population franchisse la barre, et (b) $n_{\min}(x)$ obtenu par
  \texttt{min\_ticks\_to\_admission}. Choisissez la valeur pour laquelle
  (b) est assez grand pour compter opérationnellement (plusieurs tours
  soutenus d'effort attaquant, non un seul), tout en gardant (a)
  confortablement franchissable par votre trafic réel. Il n'existe pas
  de valeur qui maximise les deux à la fois~: le corollaire~1 le prouve
  formellement, inutile d'en chercher une.
\item \textbf{Fixez $\kappa$ et $C$ en dernier, à partir de la capacité
  réelle de votre étage coûteux.} $C$ doit être votre capacité
  concurrente réelle pour l'étage coûteux (combien de dossiers il peut
  traiter à la fois sans dégradation), et $\kappa$ contrôle avec quelle
  vigueur la barre monte à mesure que cette capacité se remplit.
  Commencez à $\kappa=0{,}5$ et augmentez-le si le test d'ablation
  (section~\ref{sec:tuning} et l'attaque~5 de l'article) n'étale pas
  suffisamment les admissions concurrentes pour votre tolérance
  opérationnelle.
\item \textbf{Relancez la suite de tests adversariaux avec vos nouvelles
  constantes, non seulement le contrôle de cohérence.} Un réglage qui ne
  vérifie que l'admission des dossiers authentiques, sans relancer les
  attaques par salve et par inondation soutenue, n'a pas réellement été
  validé.
\end{enumerate}

\section{Liste de contrôle d'intégration}
\label{sec:checklist}

\begin{itemize}
\item[$\square$] Vérifié que la frontière d'admission à protéger
  correspond bien aux trois conditions de la section~\ref{sec:when} ---
  en particulier, que l'« admission » y est aujourd'hui une décision à
  signal unique, en un seul pas, non déjà limitée en débit par un autre
  moyen.
\item[$\square$] Identifié votre unité de pas réelle et votre source
  réelle d'événements de renforcement (ce qui compte concrètement, dans
  votre système, comme « une cohorte renforçant un dossier » --- un
  signalement, une réfutation, une resoumission).
\item[$\square$] Choisi un proxy d'identité de cohorte et consigné
  honnêtement sa limite. Ce prototype utilise un identifiant structurel,
  \texttt{rater\_creation\_cohort}~; quel que soit le proxy retenu,
  l'escompte d'indépendance de la porte vaut ce que vaut la résistance de
  ce proxy à l'usurpation (section~\ref{sec:limites}).
\item[$\square$] Câblé \texttt{release()} sur l'événement qui marque
  effectivement la fin du traitement d'un dossier à l'étage coûteux ---
  non sur l'événement d'admission, et non laissé sans appel
  (section~\ref{sec:pieges}, second cas).
\item[$\square$] Redérivé $x=\theta_0\rho/\gamma$ pour votre propre taux
  de renforcement des litiges authentiques, plutôt que de recopier
  telles quelles les constantes de référence (section~\ref{sec:tuning}).
\item[$\square$] Relancé les six tests de
  \texttt{tests/test\_admission\_control.py} avec vos propres constantes,
  et lu les chiffres imprimés, non seulement le statut de réussite --- le
  taux d'admission du contrôle de cohérence et le retard mesuré imposé à
  l'attaquant sont les chiffres qui indiquent si le réglage est
  réellement viable pour votre trafic.
\item[$\square$] Décidé explicitement du sort d'un dossier qui ne
  franchit jamais le seuil. Cette porte ne possède pas sa propre
  échappatoire par délai forcé~; si votre système en a besoin (une
  révision par un médiateur humain après $N$ pas de pression bloquée, par
  exemple), c'est un mécanisme distinct, superposé, non un défaut caché
  ici.
\end{itemize}

\section{Dépannage}
\label{sec:pieges}

Les deux lignes ci-dessous sont des bogues réels, découverts en
construisant et en testant adversarialement cette porte, non des modes de
défaillance hypothétiques.

\begin{center}\small
\begin{tabular}{>{\raggedright\arraybackslash}p{4.3cm}>{\raggedright\arraybackslash}p{4.3cm}>{\raggedright\arraybackslash}p{4.9cm}}
\toprule
Symptôme & Cause racine & Correctif \\
\midrule
Une inondation d'une même cohorte atteint tout de même le plafond du pas
et se fait admettre, alors même que l'escompte de répétition est
correctement configuré &
Les contributions d'une même cohorte étaient sommées \emph{avant}
déduplication à l'intérieur d'un pas, si bien qu'assez d'entrées
escomptées issues d'une seule source pouvaient tout de même atteindre le
plafond par le seul volume intra-pas &
Dédupliquer les cohortes par pas avec un simple appel à \texttt{set()}
\emph{avant} d'appliquer l'escompte et le plafond, non après \\
\addlinespace
Sous une simulation d'attaque à campagnes multiples et soutenues, le
seuil d'admission ne cesse de monter et les dossiers authentiques
ultérieurs, pourtant bien corroborés, cessent d'être admis &
\texttt{release()} n'était jamais appelé, si bien que
\texttt{\_pending\_admitted} ne fait que croître, et
\texttt{current\_threshold()} monte de façon monotone, pur artefact de
la simulation et non propriété réelle du mécanisme &
Modéliser un temps de séjour borné à l'étage coûteux et appeler
\texttt{release()} à son échéance~; ne pas assouplir l'assertion du test
pour masquer un seuil qui aurait dû redescendre \\
\bottomrule
\end{tabular}
\end{center}

La leçon générale derrière ces deux lignes~: tout ce que fait cette porte
est sensible au moment \emph{exact} où une étape de mutation s'exécute
par rapport au plafond ou au signal de charge qu'elle alimente. Chaque
fois que vous touchez à la boucle de pas, revérifiez l'ordre des
opérations, pas seulement l'arithmétique.

\section{Adapter la porte à une autre frontière d'admission}
\label{sec:adapt}

\begin{enumerate}
\item \textbf{Identifiez votre proxy de cohorte.} Tout ce qui distingue
  plausiblement des sources indépendantes convient~: âge du compte
  combiné à un identifiant de grappe de création, empreinte d'origine
  réseau, ou, plus faible mais parfois tout ce dont on dispose, simple
  identifiant utilisateur. L'escompte d'indépendance de la porte vaut ce
  que vaut la résistance de ce proxy à une usurpation renouvelée à
  chaque pas --- énoncez cette limite explicitement partout où vous
  déployez cette porte, comme le fait la section~\ref{sec:limites} ici.
\item \textbf{Identifiez votre signal de charge.} Il doit s'agir d'un
  compte réel, présentement occupé, qui se vide à mesure que des dossiers
  quittent l'étage coûteux, non d'un compteur cumulatif. Si votre étage
  coûteux n'a pas de notion naturelle d'« occupation actuelle » (par
  exemple un calcul coûteux à usage unique plutôt qu'un étage à temps de
  séjour), la couche de seuil adaptatif à la charge (règle~3 de la
  section~\ref{sec:rules}) peut ne pas se transposer~; la couche
  décroissance-plus-plafond (règles~1 et 2) se transpose, elle, seule.
\item \textbf{Redérivez la tension d'étalonnage pour votre propre
  distribution de litiges authentiques.} Ne réutilisez pas
  $\theta_0=4{,}0$ tiré du motif de renforcement d'un autre domaine~;
  relancez la procédure de la section~\ref{sec:tuning} contre votre
  propre trafic.
\item \textbf{Relancez la suite adversariale, adaptée aux formes
  d'attaque de votre domaine.} Les cinq scénarios de l'article compagnon
  (salve en un pas, inondation persistante d'une même cohorte,
  inondation persistante à cohorte usurpée, pulvérisation sur de
  nombreux dossiers, campagnes parallèles synchronisées) forment une
  liste de départ raisonnable, non une liste close --- réfléchissez à ce
  qu'un attaquant par épuisement de ressources tenterait effectivement
  en premier dans votre domaine précis.
\end{enumerate}

\section{Limites connues}
\label{sec:limites}

La porte ne ferme pas l'attaque soutenue et diversifiée en identité~;
elle la retarde, d'un facteur qui est lui-même un compromis réglable
face à la barre imposée aux dossiers authentiques (corollaire~1 de
l'article compagnon). Elle suppose que le proxy d'identité de cohorte a
un coût réel à usurper~; si un attaquant peut générer une identité
d'apparence structurellement fraîche à coût marginal nul à chaque pas,
le facteur de retard constitue toute la défense, non un filet posé sur
quelque chose de plus solide. La couche adaptative à la charge suppose
que le signal de charge est câblé sur une occupation réelle, selon le
second cas de dépannage ci-dessus --- s'y tromper laisse le comportement
de la couche indéfini par ce document. Enfin, cette porte est un contrôle
d'admission, non un mécanisme de jugement sur le contenu~: elle ne dit
rien sur la question de savoir si un dossier, une fois admis, est
correctement tranché par l'étage coûteux auquel elle donne accès.

Pour le modèle formel, les deux propositions et leurs preuves, le
corollaire sur la tension d'étalonnage et la recherche d'antériorités,
voir l'article compagnon de ce dossier.

\end{document}
